\section{Expanded reasoning: placement, joins, anchors and candidate search}
\label{sec:expand-synth}
\subsection{From stored glyph coordinates to new $x,y$ positions}
For glyph point $(u,v)$, lead $l$, shear $s$, baseline drift $d$, origin
$(p_x,p_y)$ and physical x-height $H$ in millimetres,
\FormulaNote{Purpose in this explanation: From stored glyph coordinates to new \$x,y\$ positions}{code2026}
\[
x=p_x+H(l+u+sv),\qquad y=p_y+H(v+d).
\]
Only $p_x,p_y,H$ carry millimetres; other quantities are dimensionless.
If $H=4,p_x=10,p_y=20,l=0.1,u=0.5,v=1,s=0.2,d=0.05$, the placed
point is $(13.2,24.2)$ mm. Setting shear to zero removes slant but does
not remove writer-specific stroke shape. Setting drift to zero straightens
the baseline but does not erase advance/word-gap variation.

Character advance moves the placement cursor; it is not necessarily the
glyph's bounding width. Overhanging strokes may extend beyond the advance.
Word gap is a separate measurement. Before placing a word, estimate its
advance sum and compare with remaining line width. If it does not fit, wrap
and start a new baseline. A word wider than a whole line requires a defined
overflow policy; wrapping alone cannot make it fit. Image-down coordinates
and physical-up glyph coordinates require an explicit vertical transformation
later, rather than silently changing this equation.

The reported line step is $(a_{\rm asc}-a_{\rm desc}+0.85)H$, where a
descender extent is signed negative. For ascender 1.7, descender $-0.6$,
$H=4$, the pitch is $3.15(4)=12.6$ mm. The added 0.85 is line clearance
in x-height units, not the drift persistence that happens to share the same
number. Source comments give later ruling-specific refinements; they should
not be retroactively read into every reported physical run.

\subsection{Random variation, slow drift and collision guards}
Small random translation/rotation/scale perturbations reduce mechanical
repetition. Their amplitude must remain small enough not to change character
identity. A fixed pseudorandom seed reproduces the same choices under the
same algorithm and iteration order; it does not make them naturally human.

The inspected drift update is conceptually
$v_t=0.85v_{t-1}+\eta_t$ followed by
$d_t=\operatorname{clip}(d_{t-1}+v_t,-0.25,0.25)$.
Thus persistence acts on drift \emph{velocity}, not simply on baseline
position. If noise stopped, velocity retains $0.85^k$ after $k$ glyphs;
its half-life is $\log(0.5)/\log(0.85)\simeq4.27$ glyphs.
Position is an integrated, clipped process and does not inherit a simple
stationary AR(1) variance formula. The $\pm0.25$ limit means $\pm1$ mm
at $H=4$ mm. Larger persistence produces longer correlated runs; larger
bounds allow more baseline wander. The precise values are project heuristics.

A $0.05H$ clearance is 0.2 mm at that size. It limits unwanted collisions
between separate glyph extents; legitimate cursive attachment is a different
case. Bounds-based clearance can be conservative because two overlapping
boxes need not contain overlapping strokes. Removing the guard may let
letters collide; increasing it can create unnatural loose spacing. This is
an explainable trade-off, not a demonstrated perceptual optimum.

\subsection{Automatic anchoring: piecewise formula and provenance}
Let $d$ be median stored prior distance, where low means agreement with
the character reference. The reported automatic anchoring strength is
\FormulaNote{Purpose in this explanation: Automatic anchoring: piecewise formula and provenance}{code2026}
\[
\lambda=\operatorname{clip}\left({d-0.24\over0.32-0.24},0,1\right).
\]
It equals zero for $d\le0.24$, rises linearly through 0.5 at $d=0.28$,
and saturates at one for $d\ge0.32$. The interior slope is 12.5, so a
0.01 distance change gives 0.125 in $\lambda$ before clipping. These endpoints
encode a sensitive quality transition; they are project settings, not values
derived from the cosine-distance definition.

Anchoring strength controls fallback selection policy. It is not necessarily
the interpolation weight between corresponding vertices of two strokes.
A writer's medoid, a separately stored prototype, a case-related glyph and a generic anchor
have different provenance. A generic anchor with the writer's slant applied
does not become an observed sample of that writer. If the pooled prior is
missing and the score routine returns a guarded zero, that is an absence of
evidence, not proof of a perfect glyph. Alternative policies could use
per-character confidence or a learned quality predictor, requiring labelled
quality examples and validation.

\subsection{Connectedness suppression, source compatibility and gap limits}
Measured connectedness can include accidental touching or bad character cuts.
The report suppresses automatic joins below 0.45. Unlike most exact constants,
this one has a documented \emph{intent} in source comments: example measured
values include Dylan 0.195, Owen 0.319 and Yeukita 0.417, described as mostly
print, versus writer 384 at 0.500 and 551 at 0.557, described as joined.
The chosen 0.45 lies between the two groups~\cite{code2026}.
Those comments supply a local rationale, not a newly verified dataset study
or proof that the threshold transfers to new writers. Rebuilding profiles
with different extraction can change the measurement even if the nominal
formula is unchanged.

A gap at most 0.80 x-heights limits a connector to a local relation. With
$H=4$ mm the limit is 3.2 mm. Joining additionally checks compatible sources,
eligible body strokes and endpoint geometry, so passing distance alone is
insufficient. Dots and crossbars should not become entry/exit anchors.
Height-compatible variants can be weighted more strongly without forcing
every pair to join. A shallow connector preserves readable body connections;
long straight spans can introduce false strokes or cross other glyphs.
The 0.80 limit is a geometric heuristic, not a physiological maximum.

These synthesis joins differ from the G-code writer's 0.6 mm stroke join.
The former encodes handwriting structure in x-height units; the latter
avoids a physical pen toggle between sufficiently close path endpoints.
The second can add a connector even after the first stage has finished.
Confusing them hides where a visible extra mark originated.

\subsection{Character score $A$, probability $P$ and joint score $J$}
For intended text $t$ and recognised text $\hat t$,
$A=\max(0,1-d_{\rm edit}(\hat t,t)/|t|)$ for nonempty $t$.
$P$ is the requested writer's softmax probability from the chosen classifier.
The reported combination is
\FormulaNote{Purpose in this explanation: Character score \$A\$, probability \$P\$ and joint score \$J\$}{code2026}
\[
J=\begin{cases}2AP/(A+P),&A+P>0,\\0,&A+P=0.\end{cases}
\]
This is the harmonic mean. For positive $A,P$,
$1/J=\tfrac12(1/A+1/P)$. The smaller quantity therefore has stronger
influence than under an arithmetic mean. With $A=1,P=0.1$, $J\simeq0.182$
while the arithmetic mean is 0.55. A zero in either input gives zero.
This is a project objective joining two differently interpreted scores,
not a likelihood or a percentage of physically correct handwriting.

For the report's example $A=18/19\simeq0.9474,P=0.58$, $J\simeq0.7195$.
The displayed 0.72 is rounding. Derivatives
$\partial J/\partial A=2P^2/(A+P)^2$ and
$\partial J/\partial P=2A^2/(A+P)^2$ show that when $A$ is large and
$P$ small, improving $P$ matters more. This algebra explains the preference;
it does not establish that raw writer probabilities are calibrated or
commensurate with edit accuracy.

Alternatives include a product, a weighted mean, minimum score, or maximising
one subject to a threshold on the other. Each encodes different priorities.
A product penalises two middling values more strongly; a weighted mean can
hide failure in one term; a hard constraint has a discontinuity. The harmonic
choice is defensible as a balanced preference, while its exact practical
benefit remains an empirical question.

\subsection{Full candidate and repair algorithm, with limitations}
Generate at least one candidate, using seed $s+977k$ for trial $k$;
render at controlled ink width; recognise its text; obtain writer probability;
compute $J$; retain the strictly better complete candidate. Stop at the trial
budget or the configured near-perfect score. The integer 977 spaces seeds
deterministically. Being prime does not prove statistically independent draws.
The direct joint function defaults use six tries, three repair rounds and two new
weak positions per round, with early thresholds 0.99 and 0.995. These are
search-cost limits and stopping preferences, not character-accuracy targets
guaranteed on paper.

Repair aligns the known text to find weak positions, strengthens selected
anchors, regenerates a complete trajectory and rescores it. Commit the
new anchor assignments only if the complete joint score improves. The
result's best score is then nondecreasing for deterministic scoring, but
there is no global-optimum guarantee: regenerated geometry, wrapping and
writer cues interact. A failed local change does not prove that character
cannot be improved by another choice. The removed older legibility-only routine had
different defaults and acceptance bookkeeping; it must not be described as
the joint algorithm merely because names are similar.

The same reader/classifier selecting candidates cannot provide an independent
post-selection validation. More trials can exploit model preferences and
raise $J$ without improving human legibility or physical reproduction.
The reported diagnostic candidate plots should be read as optimisation
examples. Photos, held-out read-back or separately specified geometric
measurements answer different questions. No candidate search, rendering
inference or physical reproduction was run in preparing these explanations.
